\chapter{4x10\textsuperscript{9}D.C.}\label{4x10^9-2}

\hfill\break

Ruppi un'altra lampada durante uno degli accessi di febbre.

Questa volta Diane riuscì a nasconderla al portiere. Aveva convinto il
personale delle pulizie a non fare il cambio delle lenzuola, ma a
lasciarci quelle pulite in cambio di quelle sporche ogni due giorni e
sulla porta. In questo modo le cameriere non dovevano entrare a
sistemare la stanza ed evitavamo il rischio che mi trovassero in stato
delirante. Negli ultimi sei mesi, all'ospedale locale si erano
verificati casi di dengue, colera e SDC, la sindrome da deperimento
cardiovascolare. Non volevo svegliarmi in un reparto epidemiologico
accanto a un paziente in quarantena.

«Mi preoccupa quello che potrebbe succedere quando non sono qui» mi
confessò Diane.

«Posso curarmi da solo.»

«Non se la febbre raggiunge il picco.»

«Allora è questione di fortuna e tempismo. Pensavi di andare da qualche
parte?»

«Al solito posto. Intendo in caso di emergenza, comunque. O se non posso
tornare in camera per qualche motivo.»

«Che genere di emergenza?»

Scrollò le spalle. «È un'ipotesi» ribatté con un tono che suggeriva
tutt'altro.

\separatorefiori

Non insistetti. Per migliorare la situazione non potevo far altro che
collaborare.

Stavo iniziando la seconda settimana di trattamento e mi avvicinavo alla
crisi. Il farmaco marziano si era accumulato a livelli critici nel
sangue e nei tessuti. Anche quando la febbre calava, mi sentivo
disorientato, confuso. Neanche gli effetti collaterali puramente fisici
erano divertenti. Dolori articolari. Itterizia. Eruzioni cutanee, se le
``eruzioni cutanee'' includono la sensazione di avere la pelle a
brandelli, strato dopo strato, la carne esposta quasi sanguinante come
una ferita aperta. Alcune notti dormivo per quattro o cinque ore - direi
al massimo cinque - e poi mi svegliavo in una melma di pelle morta che
Diane ripuliva dal letto chiazzato di sangue, mentre io mi spostavo su
una sedia come un artritico.

Arrivai a diffidare anche dei momenti di maggiore lucidità, che molto
spesso era puramente allucinatoria e in cui il mondo sembrava
luminosissimo e ad altissima definizione, e le parole e la memoria
parevano dentate come ingranaggi di un motore fuori controllo.

Era tremendo per me, ma forse anche peggio per Diane che doveva svuotare
la padella nei periodi in cui ero incontinente. In un certo senso mi
stava restituendo un favore. Le ero stato accanto quando anche lei era
passata per questa fase di travaglio, ma era accaduto molti anni prima.

\separatorefiori

Quasi tutte le sere dormiva accanto a me e non so come facesse. Stava
attenta a non starmi troppo vicina - certe volte, anche la sola
pressione del lenzuolo di cotone era talmente dolorosa da farmi piangere
- ma la sensazione quasi subliminale della sua presenza era
rassicurante.

Nelle notti davvero brutte, quando poteva capitare che dimenandomi
stendessi un braccio colpendola, lei si rannicchiava sul divano a fiori
accanto alle porte del balcone.

Non parlava molto dei suoi giri per Padang. Sapevo solo a grandi linee
quello che faceva: creava agganci con commissari di bordo e addetti al
carico, valutava le alternative per attraversare l'Arco. Un compito
pericoloso. Se c'era una cosa che mi faceva stare peggio degli effetti
del farmaco era vedere Diane uscire dalla porta per affrontare quella
realtà asiatica corrotta e potenzialmente violenta, senza altra
protezione che uno spray antiaggressione e il suo considerevole
coraggio.

Ciononostante, persino quel rischio intollerabile era meglio dell'essere
catturati.

Loro - e con ``loro'' intendo gli agenti dell'amministrazione Chaykin o
i loro alleati a Giacarta - erano interessati a noi per una serie di
motivi. Per il farmaco, ovviamente, e ancora di più per le varie copie
digitali degli archivi marziani che avevamo con noi. E avrebbero tanto
desiderato interrogarci sulle ultime ore di Jason: sul monologo di cui
ero stato testimone e che avevo registrato, e su tutto quello che mi
aveva raccontato riguardo alla natura degli Ipotetici e dello Spin, cose
di cui solo Jason era venuto a conoscenza.

\separatorefiori

Mi addormentai e mi svegliai, e lei era uscita.

Passai un'ora a guardare il movimento delle tende del balcone,
osservando l'angolazione che prendeva la luce solare all'estremità
visibile dell'Arco, sognando a occhi aperti le Seychelles.

Mai stati alle Seychelles? Neanch'io. Nella testa mi scorrevano le
immagini di un vecchio documentario della PBS che avevo visto una volta.
Le Seychelles sono isole tropicali, habitat di tartarughe, \emph{coco de
	mer} e una decina di varietà di uccelli rari. Geologicamente, sono tutto
ciò che resta di un antico continente che un tempo collegava l'Asia e il
Sud America, molto prima dell'evoluzione dell'uomo moderno.

I sogni, aveva detto una volta Diane, sono metafore inselvatichite. La
ragione per cui sognavo le Seychelles (me la immaginai mentre me lo
spiegava) era che mi sentivo sommerso, antico, quasi estinto.

Come un continente che stava sprofondando, inondato dalla prospettiva
della mia stessa trasformazione.

\separatorefiori

Mi addormentai di nuovo. Mi svegliai, e lei non c'era ancora.

\separatorefiori

Mi svegliai ancora solo. Era buio e sapevo che ormai era passato troppo
tempo. Cattivo segno. In precedenza, Diane era sempre tornata verso
l'imbrunire.

Mi ero agitato nel sonno. Il lenzuolo di cotone era steso sul pavimento,
a malapena visibile alla luce proveniente dalla strada e riflessa dal
soffitto intonacato. Ero infreddolito ma troppo dolorante per allungarmi
a prenderlo.

Il cielo fuori era splendidamente terso. Se stringevo i denti e
inclinavo la testa a sinistra, riuscivo a vedere qualche stella luminosa
attraverso le porte di vetro del balcone. Mi tenevo compagnia con l'idea
che magari alcune di quelle stelle, in termini assoluti, erano più
giovani di me.

Cercai di non pensare a Diane, a dove fosse e a cosa le stesse
accadendo.

E alla fine mi addormentai, percependo la luce ardente delle stelle
attraverso le palpebre, fantasmi fosforescenti che fluttuavano
nell'oscurità rossastra.

\separatorefiori

Mattina.

Perlomeno pensai che fosse mattina. C'era la luce del giorno oltre la
finestra. Qualcuno, molto probabilmente la cameriera, bussò due volte e
dal corridoio borbottò qualcosa in malese in tono seccato. E se ne andò
di nuovo.

A quel punto ero sinceramente preoccupato, anche se in quella
particolare fase del trattamento l'ansia si presentava in forma
d'irritabilità confusa. Cosa diavolo le era saltato in mente per stare
via tanto tempo e perché non era qui a tenermi la mano e spugnarmi la
fronte? L'idea che potessero averle fatto del male era indesiderata,
indimostrata, inammissibile dinanzi alla Corte.

Eppure, la bottiglia di plastica accanto al letto era vuota almeno dal
giorno prima o anche più, avevo le labbra talmente screpolate che si
stavano spaccando, e non ricordavo di aver zoppicato fino al gabinetto
di recente. Se volevo che i reni non smettessero di funzionare
completamente, dovevo andare a prendere dell'acqua dal rubinetto del
bagno.

Ma era già abbastanza difficile riuscire a non urlare solo stando
seduti.

Sollevare le gambe senza toccare il bordo del materasso era quasi
insopportabile, era come se al posto di ossa e cartilagini avessi vetri
rotti e rasoi arrugginiti.

E sebbene provassi a distrarmi pensando a qualcos'altro (le Seychelles,
il cielo), anche quel fragile conforto veniva distorto dalla lente della
febbre. Credetti di sentire la voce di Jason dietro di me, Jason che con
le mani sporche mi chiedeva di prendergli qualcosa - uno straccio o un
panno di camoscio. Uscii dal bagno con un asciugamano anziché un
bicchiere d'acqua e prima di accorgermi dell'errore ero già a metà
strada dal letto. ``Cretino. Ricomincia. Stavolta prendi la bottiglia
vuota. Riempila tutta. Riempila fino all'orlo. \emph{Segui il
	mestolino}.''

Gli passavo il panno di camoscio nel casotto del giardino dietro la
Grande Casa, dove i giardinieri tenevano gli attrezzi.

Avrà avuto circa dodici anni. All'inizio dell'estate, un paio di anni
prima dello Spin.

Tornai al presente. Sorseggiai l'acqua, assaporai il tempo e i ricordi
ripresero a fluire.

\separatorefiori

Rimasi sorpreso quando Jason propose di sistemare il tagliaerba a
benzina del giardiniere. Il giardiniere della Grande Casa era un belga
irascibile, si chiamava De Meyer, fumava Gauloises una dietro l'altra e
quando gli parlavamo si limitava a scrollare le spalle, bisbetico.
Malediceva continuamente il tagliaerba perché sputava fumo e ogni volta
si spegneva dopo qualche minuto. Perché fargli un favore? Ma era la
sfida intellettuale che affascinava Jase. Mi raccontò che era stato a
fare ricerche in rete sui motori a benzina fin dopo mezzanotte. Gli
avevano stimolato la curiosità. Aggiunse che voleva vederne uno \emph{in
	vivo}. Il fatto di non sapere cosa significasse in vivo rese la
prospettiva doppiamente interessante. Gli dissi che sarei stato felice
di aiutarlo.

In realtà mi limitai a guardare o poco più mentre posizionava il
tagliaerba su una decina di fogli del \emph{Washington Post} del giorno
prima e iniziava a esaminarlo. Questo accadeva in fondo al prato,
all'interno del capanno degli attrezzi ammuffito e isolato, dove l'aria
puzzava di nafta e benzina, fertilizzanti ed erbicidi. Sacchi di sementi
e pacciame di corteccia poggiavano incerti su scaffali di pino grezzo
tra lame azzoppate e maniglie scheggiate di attrezzi da giardino. Ci
avevano proibito di giocare nel capanno degli attrezzi. Di solito era
chiuso a doppia mandata, ma Jason ne aveva recuperato la chiave da un
gancio sulla porta del seminterrato.

Era un caldo pomeriggio di venerdì e non mi dispiaceva stare lì dentro a
guardarlo lavorare; era istruttivo e allo stesso tempo stranamente
rilassante. Per prima cosa ispezionò la macchina da terra, sdraiandovisi
accanto. Con pazienza passò le dita sulla gondola del motore,
individuando le teste delle viti e, quando fu soddisfatto, tolse le viti
e le mise da parte in ordine, e poi, accanto a queste, dopo averlo
smontato, l'alloggiamento.

E così si immerse nei meccanismi interni della macchina. In qualche modo
Jason aveva imparato da solo o intuito l'uso di un cacciavite a
cricchetto e una chiave torsiometrica. A volte le sue mosse erano
tentativi, ma mai incerti. Lavorava come un artista o un atleta: sapeva
cogliere le sfumature, era perspicace e consapevole dei propri limiti.
Aveva smontato tutte le parti raggiungibili e le aveva disposte in
bell'ordine sulle pagine annerite di grasso del \emph{Post} come
un'illustrazione anatomica, quando tutt'a un tratto la porta del capanno
si aprì stridendo e ci fece sobbalzare.

E.D. Lawton era tornato a casa presto.

«Merda» imprecai, e con questo mi aggiudicai un'occhiataccia da Lawton
padre. Rimase sulla soglia a squadrare il relitto con il suo bell'abito
grigio su misura, mentre io e Jason ci fissavamo i piedi, istintivamente
colpevoli come se fossimo stati sorpresi con una copia di
\emph{Penthouse}.

«Lo stai \emph{sistemando o distruggendo}?» chiese infine, con quel suo
tono di disprezzo misto a sdegno che era la sua firma verbale, un vezzo
che padroneggiava da così tanto tempo che ormai gli veniva naturale.

«Lo sto sistemando» rispose Jason, mansueto.

«Capisco. È \emph{tuo} il tagliaerba?»

«No, certo che no, ma pensavo che al signor De Meyer avrebbe fatto
piacere se io\ldots»

«Ma non è neanche del signor De Meyer il tagliaerba, o sbaglio? Il
signor De Meyer non possiede attrezzi suoi. A quest'ora collezionerebbe
sussidi pubblici se non lo assumessi io ogni estate. Si dà il caso che
sia il \emph{mio} tagliaerba.» E.D. attese che il silenzio si espandesse
fino a diventare quasi doloroso. Poi disse: «Hai scoperto qual è il
problema?»

«Non ancora.»

«Non ancora? Allora sarà meglio che continui.»

Jason sembrava sollevato in modo quasi soprannaturale. «Sì, ho pensato
che dopo cena mi\ldots»

«No. Non dopo cena. Lo hai smontato, e ora lo sistemi e lo rimonti. Dopo
potrai mangiare.» Poi rivolse la sua sgradita attenzione verso di me.
«Va' a casa, Tyler. Non voglio trovarti mai più qua dentro. E avresti
già dovuto saperlo.»

Sgambettai fuori alla luce abbagliante del pomeriggio battendo le
palpebre.

Non mi sorprese mai più nel capanno, ma solo perché fui attento a
evitarlo. Vi ritornai quella sera stessa, dopo le dieci, vedendo dalla
finestra della mia stanza che la luce filtrava ancora dalla fessura
sotto la porta del casotto. Presi dal frigorifero una coscia di pollo
avanzata, la avvolsi nella stagnola e mi precipitai lì con il favore
della notte. Bisbigliando, chiamai Jase che spense la luce per il tempo
necessario a farmi sgattaiolare dentro senza essere visto.

Era decorato di tatuaggi maori di grasso e olio, e aveva rimontato il
motore del tagliaerba solo a metà. Dopo che ebbe divorato qualche
boccone di pollo, gli chiesi perché ci volesse tanto.

«Potrei rimetterlo insieme in quindici minuti ma non funzionerebbe. La
parte difficile è capire esattamente cos'è che non va. E poi continuo a
peggiorare le cose. Se provo a pulire il tubo del carburante, ci entra
l'aria. Oppure si spacca la gomma. Non c'è una cosa che sia a posto. C'è
una microfrattura nell'alloggiamento del carburatore ma non so come
sistemarla. Non ho pezzi di ricambio né gli attrezzi giusti. Non sono
nemmeno sicuro di quali \emph{siano} gli attrezzi giusti.» Increspò il
viso e per un attimo pensai che si mettesse a piangere.

«Allora rinuncia» gli dissi. «Va' da E.D. e digli che ti dispiace, e te
lo fai detrarre dalla paghetta o quel che è.»

Mi fissò come se avessi detto qualcosa di nobile ma ridicolmente
semplicistico. «No, Tyler. Grazie, ma non lo faccio.»

«Perché no?»

Ma lui non rispose. Mise via la coscia di pollo e tornò ai pezzi sparsi
di quella sua follia.

Stavo per andarmene quando qualcun altro bussò pianissimo alla porta.
Jason mi fece segno di spegnere la luce. Aprì leggermente la porta e
fece entrare sua sorella.

Era evidentemente terrorizzata al pensiero che E.D. potesse trovarla lì.
Non parlava se non sussurrando. Come me, aveva portato qualcosa a Jase.
Non una coscia di pollo, ma un palmare con connessione wi-fi.

Jason si illuminò in viso quando lo vide. «Diane!» esclamò.

Lei lo zittì e mi rivolse di sghembo un sorriso nervoso.

«È solo un gadget» minimizzò, e ci fece un cenno con la testa prima di
squagliarsela.

«Lo sa benissimo che non è così» disse Jason dopo che se ne fu andata.
«Il gadget è insignificante. È la rete che è utile. Non il gadget ma la
rete.»

Nel giro di un'ora stava già consultando un gruppo di fissati della
costa occidentale che modificavano piccoli motori per gare di robot
telecomandati. Verso mezzanotte aveva trovato il modo di riparare
provvisoriamente la decina di difetti del tagliaerba. Me ne andai,
strisciai dentro casa e lo guardai dalla finestra della mia camera
mentre andava a chiamare suo padre. E.D. si trascinò svogliatamente
fuori dalla Grande Casa in pigiama e con una camicia di flanella aperta,
e rimase lì immobile con le braccia incrociate a guardare Jason che
metteva in moto il tagliaerba, il cui rumore stonava nel buio del primo
mattino. E.D. ascoltò per qualche istante, poi scrollò le spalle e fece
cenno a Jason di tornare in casa.

Jase, indugiando sulla porta, vide la luce della mia stanza attraverso
il prato e con la mano, di nascosto, mi fece un piccolo gesto.

Ovviamente le riparazioni erano provvisorie. Il giardiniere che fumava
le Gauloises si ripresentò il mercoledì seguente. Aveva ripulito circa
metà prato quando il tagliaerba si inceppò e quindi si spense una volta
per tutte. Ascoltando all'ombra degli alberi imparammo almeno una decina
di utili imprecazioni fiamminghe. Jason, la cui memoria era quasi
eidetica, si innamorò dell'espressione ``Godverdomme mijn kloten
miljaardedju!'' - che letteralmente sarebbe: ``Che Dio maledica un
milione di volte le mie palle Gesù!'', secondo ciò che Jason riuscì a
mettere insieme dal dizionario olandese-inglese della biblioteca
scolastica della Rice. Nei mesi successivi l'avrebbe usata ogni volta
che gli si rompeva un laccio delle scarpe o si bloccava un computer.

Alla fine E.D. dovette pagare per comprare un tosaerba nuovo. Al negozio
gli dissero che riparare il vecchio sarebbe costato troppo; era un
miracolo che avesse funzionato fino ad allora. Lo sentii dire da mia
madre che a sua volta l'aveva sentito da Carol Lawton. E per quanto ne
so E.D. non ne parlò mai più con Jason.

Qualche volta ci capitò di ridere di quell'episodio - ma solo mesi dopo,
quando l'amarezza era divenuta poco più di un ricordo.

\separatorefiori

Mi trascinai fino al letto pensando a Diane, che aveva fatto a suo
fratello un regalo non solo conciliatorio, come il mio, ma
effettivamente utile. E dov'era adesso? Che regalo poteva portarmi per
alleggerire il mio fardello? La sua stessa presenza sarebbe bastata.

La luce del giorno scorreva per la stanza come acqua, come un fiume
luminoso in cui ero sospeso, in cui annegavo nei minuti vuoti.

Non tutti i deliri sono vividi e frenetici. A volte sono lenti, come
rettili a sangue freddo. Guardavo le ombre strisciare come lucertole su
per i muri della stanza d'albergo. Un battito di ciglia ed era passata
un'ora. Un altro battito e calava già la notte; quando inclinai la testa
per guardare l'Arco non vi batteva più la luce solare, al suo posto
cieli scuri, nembi tropicali, lampi indistinguibili dai raptus visivi
indotti dalla febbre ma tuoni inconfondibili, un improvviso odore di
minerali bagnati proveniente dall'esterno e il rumore della pioggia che
cadeva leggera sul balcone di cemento.

E infine un altro suono: una scheda nella serratura, il cigolio dei
cardini.

«Diane» dissi (o bisbigliai, o gracchiai).

Si precipitò nella stanza. Indossava un maglione rifinito in pelle e un
cappello a tesa larga che gocciolava acqua piovana. Si fermò accanto al
letto.

«Mi spiace» disse.

«Non devi scusarti. Solo\ldots»

«Intendo che mi dispiace, Tyler, ma devi vestirti. Dobbiamo andarcene.
Subito. \emph{Ora}. C'è un taxi che ci aspetta.»

Mi ci volle del tempo per elaborare quelle informazioni. Nel frattempo
Diane iniziò a buttare roba in una valigia rigida: vestiti, documenti
sia falsi che regolari, schede di memoria, una cassetta imbottita di
boccettine e siringhe. «Non riesco ad alzarmi» tentai di dire, ma le
parole non venivano fuori giuste.

E allora un attimo dopo cominciò a vestirmi lei. Salvai un po' di
dignità sollevando le gambe senza che me lo chiedesse e stringendo i
denti invece di urlare. Poi mi sedetti e lei mi fece bere altra acqua
dalla bottiglia accanto al letto. Mi portò in bagno e qui urinai qualche
goccia vischiosa giallo canarino. «Oh cavolo,» disse, «sei completamente
disidratato.» Mi diede un altro sorso d'acqua e mi fece un'iniezione di
analgesico che mi bruciò nel braccio come veleno. «Tyler, mi dispiace
tanto!» Ma non tanto da smettere di insistere perché indossassi un
impermeabile e un cappello pesante.

Ero abbastanza vigile da sentire l'angoscia nella sua voce. «Da cosa
stiamo scappando?»

«Mettiamola così: ho avuto un incontro ravvicinato con delle persone
sgradevoli.»

«Dove andiamo?»

«Nell'entroterra. Sbrigati.»

E così ci affrettammo lungo il corridoio poco illuminato dell'hotel, giù
per una rampa di scale fino al piano terra, Diane trascinava la valigia
con la mano sinistra e mi sosteneva con la destra. Fu una bella
sgambata. Le scale, soprattutto. «Smettila di lamentarti» sussurrò un
paio di volte. E così feci. O almeno penso di averlo fatto.

Poi uscimmo nella notte. Le gocce di pioggia rimbalzavano sui
marciapiedi fangosi e sfrigolavano sul cofano di un taxi surriscaldato
vecchio di vent'anni. L'autista mi guardò con sospetto dal suo
abitacolo. Ricambiai lo sguardo. «Non è malato» lo rassicurò Diane,
facendo intendere con un gesto che avevo bevuto, lui si accigliò ma
accettò i biglietti che gli ficcò in mano.

I narcotici fecero effetto mentre eravamo nel taxi. Le strade notturne
di Padang avevano un odore cavernoso, di asfalto bagnato e pesce marcio.
Le chiazze di olio si dividevano come arcobaleni sotto le ruote
dell'auto. Lasciammo il quartiere turistico illuminato al neon ed
entrammo nel groviglio di negozi e abitazioni che si era sviluppato
attorno alla città nei trent'anni precedenti, catapecchie di fortuna che
cedevano il passo alla nuova prosperità, bulldozer parcheggiati sotto
teloni fra baracche con il tetto di lamiera. I palazzi crescevano come
funghi dai rifiuti organici dei campi abusivi. Poi passammo per la zona
industriale, con muri grigi e filo spinato, e mi riaddormentai.

Sognai non le Seychelles ma Jason. Jason e la sua passione per le reti
(``non il gadget ma la rete''), le reti che aveva creato e abitato e i
luoghi in cui quelle reti lo avevano condotto.